\subsection*{Device Program: Parallel Computation}

The kernel is designed to operate on input data, process it in parallel, and generate the output.

You repeat the code as below:

\begin{lstlisting}[language=croqtile,caption={A device helper called from tileflow}]
__device__ void kernel(int * a, int * b, int * c, int n) {
  for (int i = 0; i < n; ++i) c[i] = a[i] + b[i];
}
\end{lstlisting}

For targets that only support the \_Separated Source Programming Model\_, the code may need to be wrapped within a \texttt{\_\allowbreak{}\_\allowbreak{}cok\_\allowbreak{}\_\allowbreak{} \{\}} block. The equivalent  code is shown below:

\begin{lstlisting}[language=croqtile,caption={The cok wrapper for separate sources}]
__cok__ {
  extern "C" void kernel(int * a, int * b, int * c, int n) {
    for (int i = 0; i < n; ++i) c[i] = a[i] + b[i];
  }
} // end of __cok__
\end{lstlisting}

Here, the \texttt{extern "C"} annotation replaces the \texttt{\_\_device\_\_} keyword used in  target, as  requires C-linkage for the device functions.

Programmers must be aware that the device program follows the \emph{Single-Program-Multiple-Data (SPMD)} paradigm. However, unlike traditional  programs, the device program does not manage data movement ---whether between the host and device or across multiple storage levels within the device. Instead, the \emph{tileflow program} orchestrates these tasks in a much simpler and safer manner.
